Адаптивное обучение рассматривается как модель управления образовательным процессом, в которой контент, сложность и темп подстраиваются под состояние обучающегося. В отличие от фиксированных курсов, адаптивная система использует наблюдаемые сигналы: результаты заданий, скорость решения, ошибки и повторные попытки.

С методической точки зрения важны три блока: диагностика текущих знаний, выбор следующего шага обучения и обратная связь \autocite{kornilov2022adaptive}. На практике это означает необходимость формального представления знаний и модели принятия решений, способной работать в условиях неполной информации.
